% TOP_PROBLEM: {"id": 3364, "title": "Goldbach conjecture", "category_id": 1, "category": {"id": 1, "name": "number_theory", "display_name": "Number Theory", "description": "Properties of integers, prime numbers, Diophantine equations.", "slug": "number-theory", "order_index": 1, "created_at": "2026-07-31T15:26:25.670Z"}, "status": "open", "problem_number": "OPG-706", "set_id": 12}
% FIRST_POSTED: 2026-10-01
% ENTRY_KIND: statement_only
% UnsolvedMath ID 3364; source catalogue code OPG-706
% !TeX program = lualatex
% Definition and sources only; this file is not an LLM solution attempt.
\documentclass[11pt,a4paper]{article}
\usepackage[margin=25mm]{geometry}
\usepackage{fontspec}
\setmainfont{lmroman10-regular.otf}[BoldFont=lmroman10-bold.otf,
 ItalicFont=lmroman10-italic.otf,BoldItalicFont=lmroman10-bolditalic.otf]
\usepackage{amsmath,amssymb,mathtools}
\usepackage{unicode-math}
\setmathfont{latinmodern-math.otf}
\AtBeginDocument{\let\setminus\smallsetminus}
\usepackage{xurl,hyperref}
\hypersetup{colorlinks=true,urlcolor=blue}
\setcounter{secnumdepth}{0}
\setlength{\parindent}{0pt}
\setlength{\parskip}{5pt}
\setlength{\emergencystretch}{3em}
\begin{document}
\section*{Goldbach conjecture}\label{3364}
{\small UnsolvedMath ID 3364 · OPG-706 · Number Theory · OpenGarden}
\subsection{Definitions and mathematical statement}
Write $\mathcal P=\{2,3,5,7,\ldots\}$ for the prime positive integers.
The binary Goldbach conjecture is the universal assertion
\[
 \forall N\in2\mathbb Z_{\ge2},\qquad
                  \exists p,q\in\mathcal P\quad N=p+q.
\]
Here $2\mathbb Z_{\ge2}=\{4,6,8,\ldots\}$, and the primes $p,q$
need not be distinct. In particular $4=2+2$ and $6=3+3$ are admissible.
For an even $N>4$, any representation uses two odd primes, since
the only even prime is two.

To remove any convention about ordered pairs, define
\[
 r(N)=\#\{p\in\mathcal P:p\le N/2,\ N-p\in\mathcal P\}.
\]
The conjecture is equivalently $r(N)\ge1$ for every even $N\ge4$.
It asks for existence, without a prescribed lower bound larger than
one or an asymptotic formula for $r(N)$. A result for all sufficiently
large even integers would still require the remaining finite range
to be settled. Sums of three primes, or sums involving a semiprime,
are different representations and do not satisfy the exact target.
A counterexample must be an even integer at least four for which
all possible prime summands up to $N/2$ fail.
\subsection{Short English statement}
Is every even integer at least four a sum of two primes, which may be equal?
\subsection{Sources}
\begin{itemize}
\item[S1] ULAM.AI and the UnsolvedMath Contributors, supplied problems.json, record 3364 (OPG-706), OpenGarden collection. Catalogue identity and source transcription; CC BY 4.0.
\url{https://huggingface.co/datasets/ulamai/UnsolvedMath}.
\item[S2] Open Problem Garden, Goldbach conjecture. Original problem and discussion; the mathematical formulation below is expanded from the supplied source record.
\url{http://www.openproblemgarden.org/op/goldbach_conjecture}.
\end{itemize}
\end{document}
